\documentclass[11pt]{article}

\usepackage[T1]{fontenc}
\usepackage{amsmath,amssymb,amsthm}
\usepackage[margin=1in]{geometry}
\usepackage[hidelinks]{hyperref}
\setlength{\emergencystretch}{3em}

\newtheorem{theorem}{Theorem}
\newtheorem{definition}[theorem]{Definition}
\theoremstyle{remark}
\newtheorem{remark}[theorem]{Remark}

\title{The Moonshot Course: a Ten-Move Programme\\
with Per-Session Drift Controls}
\author{John Robert Lawson and Claude (Fable 5)}
\date{Course round, 27 July 2026}

\begin{document}
\maketitle

\begin{abstract}
This round installs the standing strategy layer
\texttt{dossier/moonshot-course.md} at John's direction, wires it into
the router and the root handoff, and registers it with the context
validator.  The course fixes a goal stack (primary: prove exactly one
Clay alternative; secondary: bank externally legible artifacts), states
a ten-move programme ordered around one organizing device, the
design-spec inversion of Definition~\ref{def:inversion}, and prescribes
pre-flight, in-flight, and post-flight drift checks for every recursive
autonomous session.  No mathematical claim is made or changed; no route
status moves.  The Clay problem is unsolved.
\end{abstract}

\section{Epistemic scope}

This is a strategy and process round.  Its only imported mathematical
fact is the peer-reviewed averaged blow-up theorem of Tao~\cite{Tao16},
used exactly as in the preceding external round: as a barrier statement
about averaging-portable budgets.  The programme below creates
obligations and sequencing, not theorems.  The course document assumes
Gate~0 (the four milestone verifications and the 2607 erratum note)
completes before the numbered moves begin, and states the fallback if a
verification fails.

\section{The organizing device}

\begin{definition}[design-spec inversion]
\label{def:inversion}
A certified exclusion budget is read in two directions at once: on the
regularity side as a constraint any closing argument may strengthen,
and on the breakdown side as a necessary condition any candidate
singularity, forced or unforced, must satisfy.  The shared object is
one machine-validated ledger of necessary conditions with exact
hypotheses and route provenance (move M1).
\end{definition}

\begin{remark}
The inversion is what makes negative rounds compound rather than
evaporate: a failed construction stage certifies a new necessary
condition, and a failed budget certifies a new survivor feature.  Both
land in the same ledger.  It also justifies working both Clay sides
without splitting the project: alternatives C and D admit a smooth
admissible force, and for a finite-time singularity the official time
decay of the force is nearly unconstraining, so the forced torus route
(ROUTE-B2) is the widest open door recorded in the possibility tree.
\end{remark}

\section{The ten moves, with route anchors}

Stated in full in the course document; recorded here for the archive.
M1 design-spec ledger (records layer; both sides).
M2 broad-shell realization, shell rung, obligation O-AVG-001a
(ROUTE-R3C kill machinery).
M3 broad-shell realization, averaged rung, obligation O-AVG-001b
(ROUTE-R3C; connects to~\cite{Tao16}).
M4 non-portable angular/triad cycle, at most three, per the standing
freeze rule (ROUTE-R3C).
M5 forced-cascade feasibility map, with the new obligation
O-FORCE-BUDGETS: re-derive the certified budgets with an admissible
force term, then split the M1 ledger into force-payable and
force-invariant constraints (ROUTE-B2).
M6 torus Pell chaining under admissible force, stagewise exact-rational
certificates (ROUTE-B2 toward alternative D).
M7 R3B gate re-triage under portability (frozen ROUTE-R3B hygiene).
M8 pincer-to-direction bridge toward the repaired 2607 antecedent
(ROUTE-R3A/R1B).
M9 HWY certificate reproduction and bridge verdict (ROUTE-B3A/B3B).
M10 standing external-review pipeline; John executes all contact.

Sequencing: Gate~0, then M1; construction spine M2, M3, M5, M6;
regularity offensive M4, M8; single-session M7, M9; M10 continuous.

\section{Drift controls}

The course prescribes checks for every session; the archive summary is:
pre-flight telescope check (task to route node to Clay alternative in
at most three links), portability check, session label, and state sync;
in-flight tripwires for schedule-conditional labelling, countermodel
scope, prose quota, unverified sources, thrash, and tree bypass;
post-flight obligation naming, mechanical validation, milestone upkeep,
handoff replacement, bank sweep, and honesty invariants.  Two failure
modes motivated the design, both observed in this repository's history:
numerology drift and the meta-work spiral.  A five-session stall with
all checks passing escalates to John rather than continuing.

\section{Process record}

\begin{enumerate}
\item Added \texttt{dossier/moonshot-course.md} (236 lines) and
registered it with the context validator (budget 245; required markers
for the goal stack, telescope check, kill criteria, and escalation).
\item Router: the course is named in the AGENTS entry points (line
budget unchanged at 54 of 55).
\item Handoff: the load-only section now opens with the strategy gate;
the removed sentence about deleted correspondence is fully covered by
the standing router rule; the 70-line budget is unchanged and all
required markers survive.
\item This report restores the round-per-batch discipline for the
course commit.
\end{enumerate}

\section{Open obligations}

\begin{itemize}
\item \textbf{Gate 0}: four milestone verifications and the 2607
erratum note, in the order queued in the audit record.
\item \textbf{O-AVG-001a/b}: shell and averaged realizations of the
broad-shell profile (moves M2, M3).
\item \textbf{O-FORCE-BUDGETS}: force-aware re-derivation of every
certified budget (move M5); until it exists, no forced-route conclusion
may cite an unforced budget.
\item \textbf{O-SPEC-LEDGER}: the M1 ledger itself, including schema
and validator coverage.
\end{itemize}

\begin{thebibliography}{9}

\bibitem{Tao16}
T. Tao,
\emph{Finite time blowup for an averaged three-dimensional
Navier--Stokes equation},
Journal of the American Mathematical Society \textbf{29} (2016),
601--674.

\end{thebibliography}

\end{document}
